\begin{thebibliography}{99}

\bibitem{padilla2025serdux}
D. E. C. Padilla, D. D. López, J. Martínez-Páez, L. Hernández, and C. C. Marroquín, 
``SERDUX-MARCIM: Maritime Cyberattack simulation using Dynamic Modeling, Compartmental Models in Epidemiology and Agent-based Modeling,'' 
\emph{International Journal of Information Security}, vol. 24, no. 3, pp. 415--430, May 2025. (Quartil Q2).

\bibitem{chernikova2023modeling}
A. Chernikova et al., 
``Modeling self-propagating malware with epidemiological models,'' 
\emph{Applied Network Science}, vol. 8, no. 1, pp. 45--62, April 2023. (Quartil Q1/Q2).

\bibitem{kato2025agent}
H. Kato, Y. Tahara, A. Ohsuga, and Y. Sei, 
``Agent-Based Analysis of Systemic Risk in Japanese Financial Institutions under Ransomware Attacks,'' 
\emph{IEICE Transactions on Information and Systems}, vol. E108-D, no. 11, pp. 2400--2412, November 2025. (Quartil Q2).

\bibitem{math2024paradigm}
J. A. Gómez, M. R. Ramos, and L. F. Santos,
``A Paradigm for Modeling Infectious Diseases: Assessing Malware Spread in Early-Stage Outbreaks,'' 
\emph{Mathematics (MDPI)}, vol. 12, no. 24, pp. 3912--3928, December 2024. (Quartil Q1).

\bibitem{ieee2025optimal}
S. R. Lahari and M. Kumar, 
``Optimal Defense Strategies Against Ransomware With Quantum Acceleration,'' 
\emph{IEEE Access}, vol. 13, pp. 18940--18955, January 2025. (Quartil Q1/Q2).

\bibitem{ieee2019spreading}
X. Yang, J. Liu, and Y. Zhang, 
``Modeling Ransomware Spreading by a Dynamic Node-Level Method,'' 
\emph{IEEE Access}, vol. 7, pp. 132456--132468, September 2019. (Quartil Q1/Q2).

\end{thebibliography}